%!TeX root=../pridetop.tex
\chapter[Chapter \thechapter]{}
	
	
\begin{figure}[t!]
\centering
\includegraphics[width=.8\linewidth]{42top}
\captionlistentry{The arrival of the Gardiners}
\end{figure}


\lettrine[lines=6,image=true]{initials/chap42h}{ad}  Elizabeth's opinion been all drawn from her own family, she could not have formed a very pleasing picture of conjugal felicity or domestic comfort. Her father, captivated by youth and beauty, and that appearance of good-humour which youth and beauty generally give, had married a woman whose weak understanding and illiberal mind had very early in their marriage put an end to all real affection for her. Respect, esteem, and confidence had vanished for ever; and all his views of domestic happiness were overthrown. But Mr Bennet was not of a disposition to seek comfort for the disappointment which his own imprudence had brought on in any of those pleasures which too often console the unfortunate for their folly or their vice. He was fond of the country and of books; and from these tastes had arisen his principal enjoyments. To his wife he was very little otherwise indebted than as her ignorance and folly had contributed to his amusement. This is not the sort of happiness which a man would in general wish to owe to his wife; but where other powers of entertainment are wanting, the true philosopher will derive benefit from such as are given.

Elizabeth, however, had never been blind to the impropriety of her father's behaviour as a husband. She had always seen it with pain; but respecting his abilities, and grateful for his affectionate treatment of herself, she endeavoured to forget what she could not overlook, and to banish from her thoughts that continual breach of conjugal obligation and decorum which, in exposing his wife to the contempt of her own children, was so highly reprehensible. But she had never felt so strongly as now the disadvantages which must attend the children of so unsuitable a marriage, nor ever been so fully aware of the evils arising from so ill-judged a direction of talents—talents which, rightly used, might at least have preserved the respectability of his daughters, even if incapable of enlarging the mind of his wife.

When Elizabeth had rejoiced over Wickham's departure, she found little other cause for satisfaction in the loss of the regiment. Their parties abroad were less varied than before; and at home she had a mother and sister, whose constant repinings at the dulness of everything around them threw a real gloom over their domestic circle; and, though Kitty might in time regain her natural degree of sense, since the disturbers of her brain were removed, her other sister, from whose disposition greater evil might be apprehended, was likely to be hardened in all her folly and assurance, by a situation of such double danger as a watering-place and a camp. Upon the whole, therefore, she found, what has been sometimes found before, that an event to which she had looked forward with impatient desire, did not, in taking place, bring all the satisfaction she had promised herself. It was consequently necessary to name some other period for the commencement of actual felicity; to have some other point on which her wishes and hopes might be fixed, and by again enjoying the pleasure of anticipation, console herself for the present, and prepare for another disappointment. Her tour to the Lakes was now the object of her happiest thoughts: it was her best consolation for all the uncomfortable hours which the discontentedness of her mother and Kitty made inevitable; and could she have included Jane in the scheme, every part of it would have been perfect.

»But it is fortunate,« thought she, »that I have something to wish for. Were the whole arrangement complete, my disappointment would be certain. But here, by carrying with me one ceaseless source of regret in my sister's absence, I may reasonably hope to have all my expectations of pleasure realized. A scheme of which every part promises delight can never be successful; and general disappointment is only warded off by the defence of some little peculiar vexation.«

When Lydia went away she promised to write very often and very minutely to her mother and Kitty; but her letters were always long expected, and always very short. Those to her mother contained little else than that they were just returned from the library, where such and such officers had attended them, and where she had seen such beautiful ornaments as made her quite wild; that she had a new gown, or a new parasol, which she would have described more fully, but was obliged to leave off in a violent hurry, as Mrs Forster called her, and they were going to the camp; and from her correspondence with her sister there was still less to be learnt, for her letters to Kitty, though rather longer, were much too full of lines under the words to be made public.

After the first fortnight or three weeks of her absence, health, good-humour, and cheerfulness began to reappear at Longbourn. Everything wore a happier aspect. The families who had been in town for the winter came back again, and summer finery and summer engagements arose. Mrs Bennet was restored to her usual querulous serenity; and by the middle of June Kitty was so much recovered as to be able to enter Meryton without tears,—an event of such happy promise as to make Elizabeth hope, that by the following Christmas she might be so tolerably reasonable as not to mention an officer above once a day, unless, by some cruel and malicious arrangement at the War Office, another regiment should be quartered in Meryton.

The time fixed for the beginning of their northern tour was now fast approaching; and a fortnight only was wanting of it, when a letter arrived from Mrs Gardiner, which at once delayed its commencement and curtailed its extent. Mr Gardiner would be prevented by business from setting out till a fortnight later in July, and must be in London again within a month; and as that left too short a period for them to go so far, and see so much as they had proposed, or at least to see it with the leisure and comfort they had built on, they were obliged to give up the Lakes, and substitute a more contracted tour; and, according to the present plan, were to go no farther northward than Derbyshire. In that county there was enough to be seen to occupy the chief of their three weeks; and to Mrs Gardiner it had a peculiarly strong attraction. The town where she had formerly passed some years of her life, and where they were now to spend a few days, was probably as great an object of her curiosity as all the celebrated beauties of Matlock, Chatsworth, Dovedale, or the Peak.

Elizabeth was excessively disappointed: she had set her heart on seeing the Lakes; and still thought there might have been time enough. But it was her business to be satisfied—and certainly her temper to be happy; and all was soon right again.

With the mention of Derbyshire, there were many ideas connected. It was impossible for her to see the word without thinking of Pemberley and its owner. »But surely,« said she, »I may enter his county with impunity, and rob it of a few petrified spars, without his perceiving me.«

The period of expectation was now doubled. Four weeks were to pass away before her uncle and aunt's arrival. But they did pass away, and Mr and Mrs Gardiner, with their four children, did at length appear at Longbourn. The children, two girls of six and eight years old, and two younger boys, were to be left under the particular care of their cousin Jane, who was the general favourite, and whose steady sense and sweetness of temper exactly adapted her for attending to them in every way—teaching them, playing with them, and loving them.

The Gardiners stayed only one night at Longbourn, and set off the next morning with Elizabeth in pursuit of novelty and amusement. One enjoyment was certain—that of suitableness as companions; a suitableness which comprehended health and temper to bear inconveniences—cheerfulness to enhance every pleasure—and affection and intelligence, which might supply it among themselves if there were disappointments abroad.

It is not the object of this work to give a description of Derbyshire, nor of any of the remarkable places through which their route thither lay—Oxford, Blenheim, Warwick, Kenilworth, Birmingham, etc., are sufficiently known. A small part of Derbyshire is all the present concern. To the little town of Lambton, the scene of Mrs Gardiner's former residence, and where she had lately learned that some acquaintance still remained, they bent their steps, after having seen all the principal wonders of the country; and within five miles of Lambton, Elizabeth found, from her aunt, that Pemberley was situated. It was not in their direct road; nor more than a mile or two out of it. In talking over their route the evening before, Mrs Gardiner expressed an inclination to see the place again. Mr Gardiner declared his willingness, and Elizabeth was applied to for her approbation.

»My love, should not you like to see a place of which you have heard so much?« said her aunt. »A place, too, with which so many of your acquaintance are connected. Wickham passed all his youth there, you know.«

Elizabeth was distressed. She felt that she had no business at Pemberley, and was obliged to assume a disinclination for seeing it. She must own that she was tired of great houses: after going over so many, she really had no pleasure in fine carpets or satin curtains.

Mrs Gardiner abused her stupidity. »If it were merely a fine house richly furnished,« said she, »I should not care about it myself; but the grounds are delightful. They have some of the finest woods in the country.«

Elizabeth said no more; but her mind could not acquiesce. The possibility of meeting Mr Darcy, while viewing the place, instantly occurred. It would be dreadful! She blushed at the very idea; and thought it would be better to speak openly to her aunt, than to run such a risk. But against this there were objections; and she finally resolved that it could be the last resource, if her private inquiries as to the absence of the family were unfavourably answered.

Accordingly, when she retired at night, she asked the chambermaid whether Pemberley were not a very fine place, what was the name of its proprietor, and, with no little alarm, whether the family were down for the summer? A most welcome negative followed the last question; and her alarms being now removed, she was at leisure to feel a great deal of curiosity to see the house herself; and when the subject was revived the next morning, and she was again applied to, could readily answer, and with a proper air of indifference, that she had not really any dislike to the scheme.

To Pemberley, therefore, they were to go.